% Generated mechanically from frozen input inputs/p01/ja/introduction.tex.
% Do not edit this generated file; edit the bound locale source instead.

% OK, start here.
%

\setcounter{chapter}{0}
\chapter{序論}



\phantomsection
\label{introduction-section-phantom}

\begin{verbatim}
Copyright (C) 2005 -- 2025 Johan de Jong
Permission is granted to copy, distribute and/or modify this
document under the terms of the GNU Free Documentation License,
Version 1.2 or any later version published by the Free Software
Foundation; with no Invariant Sections, no Front-Cover Texts,
and no Back-Cover Texts. A copy of the license is included in
the section entitled "GNU Free Documentation License".
\end{verbatim}



\section{概観}
\label{introduction-section-overview}

\noindent
Laumon と Moret-Bailly の著書（\cite{LM-B}）および Fulton ほかによる
執筆中の著作に加えて、代数スタックと、その定義に必要な代数幾何学を扱う
オープンソースの教科書があってよいと私たちは考えている。Stacks Project は、
可換代数から出発し、スキーム論と代数空間論を経て、代数スタック論の包括的な
基礎を築くことにより、この役割を果たそうとするものである。

\medskip\noindent
この資料はオンラインで読まれることを想定している。多数のページに分散した
内部参照へすばやく移動できるハイパーリンクが、その重要な特徴だからである。
ブラウザに組み込まれた PDF または DVI ビューアを使用すれば、ファイルをまたぐ
リンクも機能するはずである。

\medskip\noindent
このプロジェクトは共同作業であり、皆さんの参加を歓迎する。読んでいて見つけた
誤植や誤り、あるいは提案、追加資料、例があれば、
\href{mailto:stacks.project@gmail.com}{stacks.project@gmail.com}
まで電子メールで知らせてほしい。すべてのソースファイルを収めた tarball を
ダウンロードして展開し、make を実行すれば、DVI または PDF ビューアを用いて
ローカルに閲覧できる。LaTeX ファイルを自由に編集し、改善案を電子メールで
送ってほしい。


\section{貢献者}
\label{introduction-section-attribution}

\noindent
本書の扱う範囲は広く、ここで説明しようとする理論の発展に寄与した数学者全員を
正確かつ簡潔に挙げることは、きわめて困難である。やがてコミュニティの関心が
十分に高まり、各章に歴史的注記を書く意思のある貢献者が現れることを願っている。

\medskip\noindent
本プロジェクトへの貢献者は、本書版の標題紙および
\href{https://stacks.math.columbia.edu/tex/CONTRIBUTORS}{オンライン}
に列挙されている。ここでは、主要な貢献の一部を挙げる。
\begin{enumerate}
\item Jarod Alper は代数スタックに関する文献を論じる章を寄稿した。
『文献案内』第 \ref{guide-section-short-introductions} 節を参照されたい。
\item Bhargav Bhatt はスキームのエタール射に関する章の初稿を書いた。
『スキームのエタール射』第 \ref{etale-section-introduction} 節を参照されたい。
\item Bhargav Bhatt は
『代数補遺』第 \ref{more-algebra-section-formal-glueing} 節の初稿を書いた。
\item Kiran Kedlaya は
『降下』第 \ref{descent-section-descent-universally-injective} 節の初稿を寄稿した。
\item 以下の初稿は、
\begin{enumerate}
\item 『可換代数』第 \ref{algebra-section-oka-families} 節、
\item 『入射対象』第 \ref{injectives-section-baer} 節、および
\item 体に関する章（『体』第 \ref{fields-section-introduction} 節）
\end{enumerate}
であり、Akhil Mathew ほかの厚意により
\href{https://math.uchicago.edu/~amathew/cr.html}{The CRing Project}
から採られた。
\item Alex Perry は射影加群および Mittag--Leffler 加群に関する部分を執筆し、
『可換代数』定理 \ref{algebra-theorem-ffdescent-projectivity} の証明も寄稿した。
\item Alex Perry は Schlessinger と Rim の流儀による変形理論の章を書いた。
『形式変形理論』第 \ref{formal-defos-section-introduction} 節を参照されたい。
\item Thibaut Pugin、Zachary Maddock、Min Lee は、エタール・コホモロジーの章と
跡公式の章の基礎となった講義ノートを作成した。
『エタール・コホモロジー』第 \ref{etale-cohomology-section-introduction} 節および
『跡公式』第 \ref{trace-section-introduction} 節を参照されたい。
\item David Rydh は多数の有益なコメントを寄せ、いくつもの誤りを指摘し、
剰余ジェルブに関する部分の作成に本質的な形で協力した。また、
『空間内の亜群補遺』第 \ref{spaces-more-groupoids-section-finite} 節および
第 \ref{spaces-more-groupoids-section-etale-localize} 節の内容を発案した。
\item Burt Totaro は『例』第
\ref{examples-section-non-descending-property-projective} 節、
第 \ref{examples-section-non-effective-descent-projective} 節、および
『スタックの性質』第 \ref{stacks-properties-section-dimension} 節を寄稿した。
\item pro-エタール・コホモロジーの章（『pro-エタール・コホモロジー』第
\ref{proetale-section-introduction} 節）は、Bhargav Bhatt と Peter Scholze の
論文に基づいている。
\item Bhargav Bhatt は『例』第
\ref{examples-section-non-algebraic-hom-stack} 節および
第 \ref{examples-section-flat-not-colimit-flat-finitely-presented} 節を寄稿した。
\item Ofer Gabber は誤りを見つけて訂正を寄せたほか、
『代数多様体』補題 \ref{varieties-lemma-image-connected-component}、
『形式空間』補題 \ref{formal-spaces-lemma-completion-countably-indexed}、
『亜群補遺』第 \ref{more-groupoids-section-ind-quasi-affine} 節の内容、
『空間の性質』第 \ref{spaces-properties-section-fpqc} 節の主結果、および
『平坦性補遺』命題
\ref{flat-proposition-finite-type-injective-into-flat-mod-m} の証明を寄稿した。
\item J\'anos Koll\'ar は
『可換代数』補題 \ref{algebra-lemma-hart-serre-loc-thm} および
『局所コホモロジー』命題 \ref{local-cohomology-proposition-kollar} を寄稿した。
\item Kiran Kedlaya は
『代数補遺』第 \ref{more-algebra-section-beauville-laszlo} 節の初稿を書いた。
\item Matthew Emerton、Toby Gee、Brandon Levin は肥厚化に関するいくつかの
結果、とりわけ『スタックの射補遺』の補題
\ref{stacks-more-morphisms-lemma-reduced-diagonal}、
\ref{stacks-more-morphisms-lemma-thickening-diagonals}、および
\ref{stacks-more-morphisms-lemma-thickening-properties} を寄稿した。
\item Lena Min Ji は
『代数補遺』第 \ref{more-algebra-section-principal-radical-ideals} 節の初稿を書いた。
\item Matthew Emerton と Toby Gee は『スタックの幾何』第
\ref{stacks-geometry-section-multiplicities} 節および
第 \ref{stacks-geometry-section-dimension-of-algebraic-stacks} 節の初稿を書いた。
\end{enumerate}
